\section{Predicción Conforme}

En la práctica del aprendizaje automático y el modelado numérico, las estimaciones puntuales $\hat{y} = \hat{f}(x)$ resultan insuficientes para aplicaciones críticas, ya que no informan sobre la incertidumbre asociada a la predicción. La \textbf{Predicción Conforme} (o inferencia conforme) es un paradigma estadístico \textit{post-hoc} diseñado para envolver cualquier predictor puntual de "caja negra" y transformar sus salidas en \textbf{regiones de predicción} $\mathcal{C}(x)$ (conjuntos discretos o intervalos continuos) con garantías de cobertura estadísticamente rigurosas y no asintóticas.

\subsection{Cobertura y Eficiencia}

Dada una nueva observación de prueba $(X_{\text{test}}, Y_{\text{test}}) \in \mathcal{X} \times \mathcal{Y}$ y un nivel de error especificado por el usuario $\alpha \in (0, 1)$, se busca construir una región de predicción $\mathcal{C}(X_{\text{test}})$ que garantice la \textbf{validez} o \textbf{cobertura marginal}:

\begin{equation}
    \mathbb{P}\left( Y_{\text{test}} \in \mathcal{C}(X_{\text{test}}) \right) \ge 1 - \alpha
    \label{eq:marginal_cov}
\end{equation}

Sin embargo, exigir únicamente esta condición de cobertura es insuficiente. Un clasificador trivial que devuelva siempre el espacio completo de etiquetas $\mathcal{C}(x) = \mathcal{Y}$ (o la recta real $\mathcal{C}(x) = (-\infty, \infty)$ en regresión) satisface la cobertura con probabilidad 1, pero resulta completamente inútil desde el punto de vista práctico. Por lo tanto, la construcción de regiones de predicción óptimas requiere equilibrar dos objetivos contrapuestos:


\textbf{Cobertura:} Garantiza que la probabilidad de incluir la etiqueta verdadera sea de al menos $1 - \alpha$.

\textbf{Eficiencia e Informatividad:} Minimizar la extensión de la región predicha (el tamaño del conjunto en clasificación o la longitud del intervalo en regresión).

Además, de la \textbf{cobertura} se existen principalmente dos la \textbf{cobertura marginal}, definida por la ecuación \ref{eq:marginal_cov}, y la \textbf{cobertura condicional} \ref{eq:cond_cov}.

\begin{equation}
    \mathbb{P}\left( Y_{\text{test}} \in \mathcal{C}(X_{\text{test}}) \mid  X_{\text{test}}=x \right) \ge 1 - \alpha
    \label{eq:cond_cov}
\end{equation}

El cual exige que la garantia de $1 - \alpha$ se cumpla de forma independiente para cada valor de $\mathcal{X}$.

Esto puede generalizarse para otros tipos de coberturas derivadas de las ya mencionadas.

\begin{equation}
    \mathbb{P}\left( Y_{\text{test}} \in \mathcal{C}(X_{\text{test}}) \mid  E \right) \ge 1 - \alpha
    \label{eq:marginal_cov_E}
\end{equation}

\begin{itemize}
    \item \textbf{Marginal}: E = \{ \O \}  \text{Promedio global sobre todo el dataset de calibración y prueba} 

    \item \textbf{Condicional al objeto}: E = \{ $X_{text} = x$ \}  \text{Garantía puntual e independiente para cada coordenada $x$ específica} 

    \item \textbf{Condicional a la clase}: E = \{ $X_{text} = y$ \}  \text{Tasa de acierto garantizada por separado para cada categoría discreta $y$ (Sadinle, Vovk)-AGREGAR LA REF.}

    \item \textbf{Por grupos o atributos}: E = \{ $X_{text} \in \text{Grupo} \  g$ \}  \text{Cobertura garantizada de forma independiente para cada subgrupo demográfico o de sensores (Mondrian)-AGREGAR LA REF.}    
    
\end{itemize}


Si bien la cobertura condicional al objeto~\eqref{eq:marginal_cov_E} es el estándar ideal, puesto que asegura validez individualizada para cada entrada, se ha demostrado que garantizarla de forma exacta en muestra finita y libre de distribución es teóricamente imposible sin producir regiones de predicción triviales \cite{vock2012}. Por ello, en la práctica se recurre a relajaciones intermedias condicionando por clases, grupos o regiones espaciales.

\subsection{Regiones de Predicción Óptimas}

Para formalizar el sentido en que una región de predicción es óptima, se considera primero el caso teórico en el que la distribución conjunta $P(X,Y)$ es conocida.

\subsubsection{Caso Discreto: Clasificadores de Mínima Ambigüedad}

En clasificación multiclase con $\mathcal{Y} = \{1, \dots, K\}$, la eficiencia se mide mediante la \textbf{ambigüedad} $A(\mathcal{C})$, definida como el número esperado de etiquetas asignadas: $A(\mathcal{C}) = \mathbb{E}[|\mathcal{C}(X)|]$.

El Teorema de Optimalidad de Sadinle, Lei y Wasserman \cite{sadinle2018} demuestra que el clasificador basado en conjuntos que minimiza la ambigüedad sujeto a la restricción de cobertura objetivo está dado por los conjuntos de nivel de las probabilidades condicionales de clase $p(y|x) = \mathbb{P}(Y=y \mid X=x)$:

\begin{equation}
    \mathcal{C}^*(x) = \{y \in \mathcal{Y} : p(y|x) \ge t_\alpha\}
\end{equation}
donde $t_\alpha$ es el umbral que garantiza que la cobertura marginal sea exactamente $1 - \alpha$.




\subsubsection{Caso Continuo: Región de Longitud Mínima y Cuantiles Condicionales}
En problemas de regresión donde la respuesta es continua ($\mathcal{Y} = \mathbb{R}$), la medida de eficiencia relevante es la \textbf{longitud esperada del intervalo} (medida de Lebesgue):
\begin{equation}
    \min_{\mathcal{C}} \mathbb{E}[\text{Leb}(\mathcal{C}(X))] \quad \text{sujeto a} \quad \mathbb{P}(Y \in \mathcal{C}(X)) \ge 1 - \alpha
\end{equation}

Bajo condiciones de regularidad, la región de longitud mínima corresponde a las \textbf{Regiones de Alta Densidad} (HDR, \textit{High Density Regions}) definidas sobre la densidad condicional $f(y|x)$:
\begin{equation}
    \mathcal{C}^*(x) = \{y \in \mathbb{R} : f(y|x) \ge t_\alpha\}
\end{equation}

Cuando se restringe la salida a un único intervalo continuo con colas balanceadas $\frac{\alpha}{2}$, la solución óptima está delimitada por las funciones de \textbf{cuantiles condicionales verdaderos} $q_{\alpha / 2}(x)$ y $q_{1-\alpha / 2}(x)$:

\begin{equation}
    \mathcal{C}^*(x) = \left[ q_{\alpha/2}(x), \, q_{1-\alpha/2}(x) \right]
\end{equation}

\subsection{Inferencia Conforme en Muestra Finita}

En la práctica, las distribuciones verdaderas $p(y|x)$ o $q_\tau(x)$ no son conocidas y deben ser estimadas mediante modelos de aprendizaje automático (como redes neuronales, PINNs, etc). Evaluar directamente las regiones estimadas $\hat{\mathcal{C}}(x)$ sin calibración genera intervalos sin garantías estadísticas en muestra finita.

La \textbf{Predicción Conforme} resuelve este problema actuando como un envoltorio (\textit{wrapper}) no paramétrico que utiliza una muestra de calibración independiente para ajustar los umbrales de corte mediante \textbf{puntajes de no conformidad} $s(x, y)$. Bajo el único supuesto de \textbf{intercambiabilidad} de las observaciones, la inferencia conforme garantiza validez exacta en muestra finita:

\begin{equation}
    1-\alpha \le \mathbb{P}\left(Y_{\text{test}} \in \mathcal{C}(X_{\text{test}})\right) < 1-\alpha + \frac{1}{n+1}
    \label{eq:marg}
\end{equation}

donde $n$ es el número de muestras en el conjunto de calibración.
%-------------------------------------------------------------------------------------



Este no produce una predicción puntual $\hat{y}$ sino que produce una \textbf{región de predicción} $\mathcal{C}(x)$ (un conjunto o un intervalo) estadísticamente rigurosos.
De manera crucial, estos conjuntos son válidos en un sentido libre de distribución: poseen garantías explícitas y no asintóticas incluso sin asumir una distribución de los datos ni supuestos sobre el modelo.


Dada una muestra de prueba $X_{\text{test}}$, se dispone de un modelo previamente entrenado $\hat{f}(\cdot)$ que produce una estimación puntual $\hat{y} = \hat{f}(X_{\text{test}})$. El marco de la Predicción Conforme trata a este predictor como una "caja negra" (\textit{black-box}) totalmente agnóstica.

En la ecuación \ref{eq:marg} se da una garantia marginal es decir la probabilidad en promedio sobre la aleatoriedad de todo el conjunto de datos de calibración y de la nueva muestra de prueba $(X_{\text{test}}, Y_{\text{test}})$. Esto garantiza una tasa de acierto global del $(1-\alpha)\%$ a largo plazo.


Para calibrar el modelo y obtener esta garantía, se incorpora una etapa adicional al procedimiento tradicional de entrenamiento y evaluación: el \textbf{paso de calibración}. 






Este procedimiento no requiere asumir una distribución de 
probabilidad específica para los datos, 
sino únicamente el cumplimiento de la propiedad de \textbf{intercambiabilidad}.


\begin{definicion}[\textbf{Intercambiabilidad}]
Sean $Z_i = (X_i, Y_i)$ para $i = 1, \dots, n$ las muestras del conjunto de calibración y $Z_{n+1} = (X_{\text{test}}, Y_{\text{test}})$ una nueva muestra de prueba. Se dice que la secuencia $(Z_1, \dots, Z_{n+1})$ es \textbf{intercambiable} si para toda permutación $\pi$ del conjunto de índices $\{1, \dots, n+1\}$ se cumple que:
\begin{equation}
    (Z_1, \dots, Z_{n+1}) \stackrel{d}{=} (Z_{\pi(1)}, \dots, Z_{\pi(n+1)}),
\end{equation}
es decir, la distribución conjunta de las variables aleatorias permanece inalterada frente a cualquier reordenamiento de las observaciones.
\end{definicion}


Cabe destacar que la intercambiabilidad es una hipótesis más general que 
la condición clásica de variables independientes e idénticamente 
distribuidas (i.i.d.): toda secuencia i.i.d. es intercambiable, 
pero esta propiedad también se satisface bajo ciertas dependencias 
simétricas entre los datos.

La relevancia práctica de este supuesto radica en que, 
al evaluar cada par $(X_i, Y_i)$ mediante una función de discrepancia 
fija $s(X_i, Y_i)$ basada en el modelo previamente entrenado $\hat{f}$, 
las puntuaciones resultantes $(s_1, \dots, s_n, s_{n+1})$ preservan la 
intercambiabilidad. 
En consecuencia, la puntuación desconocida del ejemplo de prueba 
$s_{n+1}$ tiene exactamente la misma probabilidad de ocupar cualquier 
rango al ordenarse junto con las $n$ puntuaciones de calibración. 
Esto permite acotar el error esperado calculando simplemente el cuantil 
empírico $\hat{q}$ correspondiente al nivel $\lceil(n+1)(1-\alpha)\rceil$ 
sobre el conjunto de calibración. 
Este procedimiento general se formaliza 
en el Algoritmo~\ref{alg:conformal_prediction}.

\begin{algorithm}[H]
\caption{Proceso de Calibración en \textit{Marco general - Predicción Conforme}}
\label{alg:conformal_prediction}
\begin{algorithmic}[1] % El [1] numera cada paso automáticamente
    \STATE \textbf{Identificar}  y extraer del modelo previamente entrenado $\hat{f}$ una estimación inicial de confianza o incertidumbre (aún sin garantías estadísticas).
    
    \STATE \textbf{Definir} la función de puntuación $s(x,y) \in \mathbb{R}$, también llamada función de no conformidad (los puntajes más grandes definen un peor ajuste entre $x$ e $y$):
    \begin{itemize}
        \item $s(x,y)$ pequeño $\rightarrow$ buen ajuste.
        \item $s(x,y)$ grande $\rightarrow$ mal ajuste.
    \end{itemize}

    \STATE \textbf{Computar} $\hat{q}$ como el $\lceil(n+1)(1-\alpha)\rceil/n$-cuantil empírico de la función de puntuación evaluada en el conjunto de calibración:
    \begin{equation*}
        \hat{q} = \text{quantile}\left(s_1, \dots, s_n; \frac{\lceil(n+1)(1-\alpha)\rceil}{n}\right)
    \end{equation*}
    donde $s_i = s(X_i, Y_i)$ para $i=1, \dots, n$. Esto fija el tamaño del conjunto de predicción de forma que se garantiza una cobertura de $1-\alpha$.

    \STATE \textbf{Construir} la región o conjunto de predicción para un nuevo ejemplo $X_{\text{test}}$ utilizando el umbral $\hat{q}$:
    \begin{equation*}
        C(X_{\text{test}}) = \{ y : s(X_{\text{test}}, y) \le \hat{q} \}
    \end{equation*}
    Es decir, el conjunto incluye todos los valores de $y$ que poseen bajo puntaje según el umbral calibrado.
\end{algorithmic}
\end{algorithm}

El Algoritmo \ref{alg:conformal_prediction} define el \textit{Marco general} sobre el cual se construyen las distintas técnicas de \textit{Conformal Prediction}.

Este flujo establece una plantilla estándar de calibración: cualquier variante metodológica posterior se reduce, esencialmente, a proponer una definición particular para la función de puntuación de no conformidad $s(x,y)$ en el paso 2, manteniendo el resto del proceso de cuantificación de incertidumbre inalterado con distintas formas de calcular el percentil.

\begin{figure}[H]
    \centering
    \resizebox{1.04\textwidth}{!}{% <-- Cambia 0.95 para agrandar (ej. 1.0) o achicar
        %-------------------------------------------------------------------------%
% ARCHIVO: img_tex/diagrama.tex
% Arreglado con desactivación local para babel (spanish)
%-------------------------------------------------------------------------%
\shorthandoff{>} % Apaga el conflicto del español justo antes de empezar

\begin{tikzpicture}[
    scale=0.85, every node/.style={transform shape},
    % Estilos de los bloques centrales
    bloque/.style={rectangle, draw=gray!60, fill=blue!5, rounded corners=6pt, minimum width=3.5cm, minimum height=1.2cm, align=center, font=\small},
    % Estilo para el texto de las ecuaciones
    textoEq/.style={align=left, font=\small}
]

% 1. FONDOS DE COLORES (Izquierda Verde, Derecha Rosa)
\fill[green!8] (-9, -5.4) rectangle (0, 5.5); % Subí un poquito el fondo (de 5 a 5.5) para que tape la nueva ecuación
\fill[red!5] (0, -5.4) rectangle (9, 5.5);
\draw[gray!30, dashed] (0, -5) -- (0, 5.5); 

% 2. BLOQUES CENTRALES Y FLECHAS
% Agregamos la ecuación arriba de Data en (0, 4.9)
\node[font=\small] (dataset) at (0, 4.99) {$(X_i, Y_i)_{i=1,2,\dots, n}$};

\node[bloque] (data) at (0, 4) {Datos};
\node[bloque] (sistema) at (0, 1.5) {Sistema};
\node[bloque] (cp) at (0, -1.5) {CP \\ $\mathbb{P}(y_{\text{test}} \in \mathcal{C}(X_{\text{test}})) \ge 1-\alpha$};

\node[textoEq] at (0, -4.9) {$y_{test} \mathcal \in {C}(x_{test}) $};

% Flecha delgada desde la nueva ecuación hacia el bloque Data
\draw[->, gray!80, thin] (dataset) -- (data);

\draw[->, gray!80, thick] (data) -- (sistema);
\draw[->, gray!80, thick] (-1, 0.9) -- node[left, black] {$S_i$} (-1, -0.9);
\draw[->, gray!80, thick] (1, 0.9) -- node[right, black] {$\hat{y}_i$} (1, -0.9);
\draw[->, gray!80, thick] (cp) -- (0, -3.5);

% Parámetro alpha apuntando a CP
\draw[->, gray!80, thick] (-2.7, -1.5) node[left, black, font=\large] {$\alpha$} -- (cp);

% 3. LADO IZQUIERDO (Clasificación - Verde)
\node[textoEq] at (-6.5, 1.5) {
    $x_1 \rightarrow \{\hat{y}_2\}$ \\
    $x_2 \rightarrow \{\hat{y}_K\}$ \\
    $\vdots$ \\
    $x_N \rightarrow \{\hat{y}_1\}$
};

\node[textoEq] at (-6.5, -2) {
    $x_1 \rightarrow \{\hat{y}_2, \hat{y}_4, \hat{y}_6\}$ \\
    $x_2 \rightarrow \{\hat{y}_K, \hat{y}_5, \hat{y}_7, \hat{y}_8\}$ \\
    $\vdots$ \\
    $x_N \rightarrow \{\hat{y}_1\}$
};

\node[textoEq] at (-3, -4.2) {$\mathcal{C}(x_{i}) = \{\hat{y}_j, \hat{y}_n, \textcolor{green!60!black}{\hat{y}_{i}}, \hat{y}_k\}$};


% 4. LADO DERECHO (Regresión - Rosa)

% Gráfico Superior (Predicción puntual)
\begin{scope}[shift={(3.5, 0.5)}]
    \draw[->, thick] (0,0) -- (3,0);
    \draw[->, thick] (0,0) -- (0,2);
    % Curva azul de predicción
    \draw[blue!70!cyan, ultra thick] (0.2, 0.4) .. controls (1, 1.5) and (2, 1) .. (2.8, 1.8) node[right, black] {$\hat{y}$};
\end{scope}

% Gráfico Inferior (Intervalos de confianza / bandas)
\begin{scope}[shift={(3.5, -3.2)}]
    \draw[->, thick] (0,0) -- (3,0);
    \draw[->, thick] (0,0) -- (0,2);
    
    % Banda superior (azul claro)
    \draw[blue!30, thick] (0.2, 0.9) .. controls (1, 2) and (2, 1.5) .. (2.8, 2.3);
    % Banda inferior (azul claro)
    \draw[blue!30, thick] (0.2, 0.1) .. controls (1, 1.2) and (2, 0.7) .. (2.8, 1.3);
    % Curva central (gris gruesa)
    \draw[gray!50, line width=3pt] (0.2, 0.5) .. controls (1, 1.6) and (2, 1.1) .. (2.8, 1.8) node[right, black] {$\hat{y}$};
    
    % Etiquetas de las bandas I+ e I-
    \draw[red!70, thick] (2.4, 1.55) -- node[right, font=\tiny] {$I^+$} (2.4, 2.0);
    \draw[red!70, thick] (0.6, 0.5) -- node[right, font=\tiny] {$I^-$} (0.6, 0.85);
\end{scope}

\node[textoEq] at (3.5, -4.2) {$\mathcal{C}(x) = [\hat{y}(x) - I(x), \hat{y}(x) + I(x)]$};

\end{tikzpicture}

\shorthandon{>} % Lo vuelve a encender para que el resto del texto de tu informe funcione bien
    }
    \caption{Esquema general de Conformal Prediction para clasificación y regresión.}
    \label{fig:conformal_prediction}
\end{figure}

El Algoritmo~\ref{alg:conformal_prediction} está pensado usualmente para el caso de clasificación. Este flujo también se representa en la Figura~\ref{fig:conformal_prediction}, donde nos interesa analizar el lado izquierdo.

Para aplicarlo, se evalúan los datos en el modelo $\hat{f}$ con el fin de extraer sus predicciones o salidas (como la salida de \textit{softmax}) y calcular así las puntuaciones mediante la función de no conformidad seleccionada.

En el caso de clasificación, esta función puede definirse como $s(x_i,y_i) = \mathbf{P}(y_i=0 \mid x_i) = 1 - \mathbf{P}(y_i=1 \mid x_i)$ (con probabilidad denotada $\hat{\pi}(y \mid x)$), mientras que para la calibración en regresión puede utilizarse $s(x_i,y_i) = |y_i - \hat{y}|$. En ambos casos, la función de no conformidad se emplea para decidir qué muestras se incluyen dentro del conjunto de predicción.

\section{Split conformal predictión}

El esquema presentado en el Algoritmo~\ref{alg:conformal_prediction} corresponde a la variante 
conocida en la literatura como \textbf{Predicción Conforme Dividida} (\textit{Split Conformal Prediction}, 
SCP) o inductiva. Para comprender su relevancia práctica, es necesario contrastarla 
con la formulación original del marco, 
denominada \textbf{Predicción Conforme Completa} (\textit{Full Conformal Prediction}).

En la Predicción Conforme Completa no se reserva un conjunto de calibración separado, 
sino que se utilizan todos los datos disponibles para ajustar el modelo. 
Dado un conjunto de $n$ observaciones $\mathcal{D}_n = \{(X_i, Y_i)\}_{i=1}^n$ y 
una nueva entrada de prueba $X_{n+1}$, el método debe evaluar cada valor posible 
$y \in \mathcal{Y}$ como candidato para la etiqueta verdadera $Y_{n+1}$.  
Para cada candidato $y$, se construye un conjunto de datos aumentado $\mathcal{D}_n \cup \{(X_{n+1}, y)\}$, 
se reentrena el modelo predictivo $\hat{f}_y$ desde cero sobre ese conjunto y se recalculan todas 
las puntuaciones de no conformidad. Si bien este procedimiento aprovecha la totalidad 
de los datos para el entrenamiento, resulta computacionalmente prohibitivo en la práctica: 
en problemas de regresión continua exige ajustar el modelo infinitas veces por cada muestra 
de prueba, y en aprendizaje profundo vuelve inviable el costo de reentrenar redes 
neuronales complejas o redes guiadas por la física (PINNs) en tiempo de inferencia.

En este caso se parte de un set de datos $(X_i, Y_i)_{i=1,2, \dots n}$ el modelo se entrena una vez con los datos $(X_i, Y_i)_{i=1,2, \dots n/2}$ y se calibra con $(X_i,Y_i)_{i=n/2 + 1, \dots , n}$
Se considera el problema de regresión donde $\mathcal{X} \in \mathbb{R}$.



\begin{algorithm}[H]
\caption{SCP caso especial - regresión}
\label{alg:conformal_prediction_esp}
    \begin{algorithmic}[1] % Usa [1] para numerar cada línea, o déjalo vacío [] si no quieres números
        \STATE \textbf{Entrenamiento:} Utilizar los datos $(X_i, Y_i)$ para $i = 1, \dots, \frac{n}{2}$ para ajustar un modelo predictivo $\hat{f} : \mathcal{X} \to \mathbb{R}$.

        \STATE \textbf{Cálculo de residuos:} Para $i = \frac{n}{2} + 1, \dots, n$, calcular el residuo absoluto $S_i = |Y_i - \hat{f}(X_i)|$.

        \STATE \textbf{Cálculo del cuantil:} Ordenar los residuos $S_{\frac{n}{2}+1}, \dots, S_n$ en orden creciente, y definir $\hat{q}$ como el elemento de valor $\lceil(1 - \alpha)(\frac{n}{2} + 1)\rceil$ en la lista ordenada.
        Entonces, $\hat{q} = \text{quantile}\left(s_1, \dots, s_n; \frac{\lceil(n+1)(1-\alpha)\rceil}{n}\right)$.

        \STATE \textbf{Inferencia:} Devolver el intervalo de predicción $\mathcal{C}(X_{n+1}) = [\hat{f}(X_{n+1}) - \hat{q}, \hat{f}(X_{n+1}) + \hat{q}]$.
    \end{algorithmic}
\end{algorithm}

El algortimo \ref{alg:conformal_prediction_esp} a diferencia del algoritmo \ref{alg:conformal_prediction} el conjunto de prediccion devuelto es un intervalo $\mathcal{C}(X_{n+1}) = [\hat{f}(X_{n+1}) - \hat{q}, \hat{f}(X_{n+1}) + \hat{q}]$ el cual representa la incerteza alrededor de la prediccion $\hat{f}(X_{n+1})$ para el target $Y_{n+1}$.

El valor de $\hat{q}$ es una constante que se calcula con los datos de calibración $(X_i,Y_i)_{i=n/2, \dots, n}$ y en esos datos se eleigio $\hat{q}$ de forma que cumpla que los $Y_i$ esten contenidos en el intervalo una fracción de $1-\alpha$. siendo los datos iid luego cualquier nueva prediccion va a cumplir la garantia marginal.


\begin{algorithm}[H]
\caption{SCP caso general}
\label{alg:conformal_prediction_gen}
    \begin{algorithmic}[1]
    \STATE \textbf{Construcción de la función de puntuación:} Utilizar los datos $(X_i, Y_i)$ para $i = 1, \dots, n/2$ para construir una \textbf{score de no conformida} $s : \mathcal{X} \times \mathcal{Y} \to \mathbb{R}$.

    \STATE \textbf{Cálculo de puntuaciones:} Para $i = n/2 + 1, \dots, n$, calcular la puntuación $S_i = s(X_i, Y_i)$ para cada dato en el conjunto de calibración.

    \STATE \textbf{Cálculo del cuantil:} Ordenar las puntuaciones $S_{n/2+1}, \dots, S_n$ en orden creciente, y definir $\hat{q}$ como el cuantil empirico de nivel $\lceil(1 - \alpha)(\frac{n}{2} + 1)\rceil$ en la lista ordenada.

    \STATE \textbf{Inferencia:} Devolver el conjunto de predicción $\mathcal{C}(X_{n+1}) = \{y \in \mathcal{Y} : s(X_{n+1}, y) \leq \hat{q}\}$.
    \end{algorithmic}
\end{algorithm}

La intuición de $s(x, y)$ es que mide qué tan ``inusual'' es el punto $(x, y)$ basándose en un modelo ajustado con los datos de esta primera mitad.

En el Algoritmo~\ref{alg:conformal_prediction_gen} vemos que este método no es solo de regresión, sino que se ajusta a ambos casos (clasificación y regresión). En general, todas las variantes que veremos en las siguientes secciones también pueden aplicarse a ambos problemas de calibración.

\section{Puntuaciones de no conformidad}

Las \textbf{puntuaciones de no conformidad} constituyen el elemento central del marco de la Predicción Conforme. Formalmente, se definen mediante una función $s(x, y)$ diseñada de modo tal que valores más altos indiquen una mayor discrepancia entre el par $(x, y)$ y el comportamiento aprendido por el modelo en los datos de entrenamiento. En términos prácticos, esta función actúa como una métrica que cuantifica el error o la falta de ajuste del predictor frente a una observación dada.

Una propiedad destacable de este enfoque es su flexibilidad: es posible emplear prácticamente cualquier función de puntuación, ya que, siempre que se preserve la simetría respecto al orden de los datos, el procedimiento garantiza matemáticamente el nivel de cobertura marginal $1-\alpha$, independientemente de cuán preciso o deficiente sea el modelo base. Sin embargo, aunque la validez estadística está asegurada por construcción, la elección específica de $s(x, y)$ resulta determinante en la práctica. Una función de no conformidad bien diseñada permite aprovechar la información del predictor para generar conjuntos o intervalos de predicción más estrechos y precisos. En la literatura del área, esta cualidad de minimizar la amplitud o el tamaño del conjunto devuelto sin perder el nivel de cobertura exigido se denomina \textbf{eficiencia}.

\subsection{Funciones de puntuación en regresión}

En problemas de regresión, la cobertura marginal por sí sola resulta insuficiente, ya que un intervalo excesivamente amplio e inflexible podría cumplirla de forma trivial perdiendo toda utilidad práctica. Idealmente, un intervalo de predicción debe reflejar la incertidumbre local de los datos, ensanchándose en zonas de alta varianza (heterocedasticidad) y estrechándose en regiones más estables. Según cómo se defina la puntuación de no conformidad, se obtienen distintos grados de adaptabilidad.

La formulación más directa es la \textbf{puntuación por residuo absoluto}, definida como 

\begin{equation}
    s(x, y) = |y - \hat{f}(x)|.
\end{equation}

Al depender únicamente del error puntual sin ponderar la dispersión local, este criterio produce bandas de confianza de \textbf{ancho constante} alrededor de la predicción, dadas por $[\hat{f}(X_{n+1}) - \hat{q}, \, \hat{f}(X_{n+1}) + \hat{q}]$.

Para superar esta rigidez cuando la varianza del ruido no es constante, se recurre a la \textbf{puntuación por residuo escalado}, que normaliza el error mediante una estimación local de la dispersión $\hat{\sigma}(x)$:
\begin{equation}
    s(x, y) = \frac{|y - \hat{f}(x)|}{\hat{\sigma}(x)}.
\end{equation}
De este modo, el intervalo de predicción resultante adapta su amplitud según la incertidumbre estimada en cada punto:
\begin{equation}
    \mathcal{C}(X_{n+1}) = \hat{f}(X_{n+1}) \pm \hat{q}\,\hat{\sigma}(X_{n+1}).
\end{equation}

Una alternativa aún más flexible es la \textbf{Regresión Cuantílica Conforme} (\textit{Conformalized Quantile Regression}, CQR), que parte de estimaciones iniciales de cuantiles condicionales $[\hat{\tau}(x; \alpha/2), \, \hat{\tau}(x; 1 - \alpha/2)]$ para medir la distancia con signo del valor real $y$ respecto a dicho intervalo:
\begin{equation}
    s(x, y) = \max \left\{ \hat{\tau}(x; \alpha/2) - y, \, y - \hat{\tau}(x; 1 - \alpha/2) \right\}.
\end{equation}
Un valor positivo de $s(x,y)$ indica que el dato real quedó por fuera del intervalo estimado, mientras que un valor negativo señala que cayó en su interior. Tras calibrar el cuantil empírico $\hat{q}$ sobre estas puntuaciones, el conjunto de predicción final se obtiene ajustando ambos extremos del intervalo inicial:
\begin{equation}
    \mathcal{C}(X_{n+1}) = \left[ \hat{\tau}(X_{n+1}; \alpha/2) - \hat{q}, \, \hat{\tau}(X_{n+1}; 1 - \alpha/2) + \hat{q} \right].
\end{equation}
Este enfoque no paramétrico permite construir \textbf{intervalos asimétricos y adaptativos} a la distribución local de los datos sin asumir simetría en los residuos.

\subsection{Funciones de puntuación en clasificación}

En problemas de clasificación, donde la variable respuesta toma valores en un conjunto discreto de etiquetas $\mathcal{Y} = \{1, \dots, K\}$, el modelo suele proporcionar una distribución de probabilidad condicional estimada $\hat{\pi}(y \mid x)$. A partir de ella, una elección habitual es la \textbf{puntuación de alta probabilidad}, definida como el negativo de la probabilidad asignada a la clase verdadera:
\begin{equation}
    s(x, y) = -\hat{\pi}(y \mid x).
\end{equation}
El signo negativo asegura la convención fundamental del marco conforme: las observaciones más improbables reciben puntuaciones más altas (más cercanas a cero).

Una vez calculado el umbral de calibración $\hat{q}$, el conjunto de predicción $\mathcal{C}(X_{n+1})$ para una nueva entrada incorpora todas las etiquetas cuya probabilidad estimada supere dicho nivel de corte:
\begin{equation}
    \mathcal{C}(X_{n+1}) = \{y \in \mathcal{Y} : \hat{\pi}(y \mid X_{n+1}) \ge -\hat{q}\}.
\end{equation}
De forma equivalente, el valor $t^* = -\hat{q}$ opera como un umbral de decisión calibrado sobre las salidas probabilísticas del clasificador.

\begin{figure}[H]
    \centering
    \resizebox{1.0\textwidth}{!}{%-------------------------------------------------------------------------%
% ARCHIVO: img_tex/diagrama_completo.tex
% Versión corregida de ALTA FIDELIDAD - Método Hercule Poirot
%-------------------------------------------------------------------------%
\shorthandoff{>} % Evita conflictos con babel spanish

% 1. DEFINICIÓN DE DATOS GLOBALES
\newcommand{\puntosRegresion}{
    0.4/2.6, 0.4/3.0, 0.5/2.9, 0.5/2.4, 0.6/2.7, 0.7/2.5, 0.7/2.2, 0.8/2.7, 
    0.9/2.4, 1.0/2.2, 1.1/2.0, 1.2/1.5, 1.3/2.3, 1.3/1.9, 1.4/1.6, 1.5/2.1, 
    1.6/1.8, 1.7/1.5, 1.8/1.7, 1.8/2.2, 1.9/1.9, 2.0/1.5, 2.1/1.6, 2.2/2.1, 
    2.3/1.8, 2.3/2.5, 2.4/2.4, 2.4/2.7, 2.5/2.6, 2.5/2.9, 2.6/2.0, 2.7/1.8, 
    2.7/2.2, 2.8/2.4, 2.8/1.6, 2.9/1.3, 2.9/2.4, 3.0/1.4, 3.0/2.1, 3.1/1.6, 
    3.2/1.5, 3.3/1.1, 3.3/2.3, 3.4/1.3, 3.4/1.8, 3.5/0.8, 3.6/0.9, 3.6/1.2, 
    3.7/1.1, 3.8/0.6, 3.9/0.4, 4.0/0.5
}

% Trayectoria de la línea de predicción central
\newcommand{\lineaPrediccion}{
    (0.5, 2.8) .. controls (1.0, 2.0) and (1.5, 1.7) .. (2.3, 1.7)
               .. controls (3.0, 1.7) and (3.5, 1.2) .. (4.0, 0.7)
}

\begin{tikzpicture}[
    scale=0.8, every node/.style={transform shape},
    eje/.style={thick, draw=black!80},
    punto/.style={circle, fill=black!50, inner sep=0.4pt},
    barra_in/.style={draw=black!70, fill=yellow!60, line width=0.5pt},
    barra_out/.style={draw=black!70, fill=gray!40, line width=0.5pt},
    banda/.style={fill=blue!20!cyan!20, opacity=0.6}
]

% =========================================================================
% FILA SUPERIOR: REGRESIÓN
% =========================================================================

% --- GRÁFICO 1: Residual Score (Ancho Constante) ---
\begin{scope}[shift={(0,5)}]
    \node[font=\large\bfseries, text=black!70] at (2.2, 3.8) {Score residual };
    \draw[eje] (0,0) -- (4.5,0) node[below left] {$\mathcal{X}$};
    \draw[eje] (0,0) -- (0,3.5) node[above left] {$\mathcal{Y}$};
    
    \fill[banda] 
        (0.5, 3.3) .. controls (1.0, 2.5) and (1.5, 2.2) .. (2.3, 2.2) 
                   .. controls (3.0, 2.2) and (3.5, 1.7) .. (4.0, 1.2) --
        (4.0, 0.2) .. controls (3.5, 0.7) and (3.0, 1.2) .. (2.3, 1.2)
                   .. controls (1.5, 1.2) and (1.0, 1.5) .. (0.5, 2.3) -- cycle;
                   
    \draw[gray!60, thick] \lineaPrediccion;
    
    % Expansión controlada y limpia mediante una macro temporal ejecutada al instante
    \edef\ejecutarLoopUno{\noexpand\foreach \noexpand\x/\noexpand\y in {\puntosRegresion}{\noexpand\node[punto] at (\noexpand\x,\noexpand\y) {};}}
    \ejecutarLoopUno
\end{scope}

% --- GRÁFICO 2: Scaled Residual Score (Efecto Globo Central) ---
\begin{scope}[shift={(5.5,5)}]
    \node[font=\large\bfseries, text=black!70] at (2.2, 3.8) {Residuo scalado};
    \draw[eje] (0,0) -- (4.5,0) node[below left] {$\mathcal{X}$};
    \draw[eje] (0,0) -- (0,3.5) node[above left] {$\mathcal{Y}$};

    \fill[banda] 
        (0.5, 2.92) .. controls (0.79, 2.58) and (0.92, 2.42) .. (1.25, 2.22)
                    .. controls (2.12, 1.82) and (2, 2.58) .. (2.8, 2.58)
                    .. controls (3.48, 2.32) and (3.01, 1.69) .. (4.03, 1.1) --
        (4.03, 0.4) .. controls (3.5, 0.82) and (3.25, 1.18) .. (2.25, 1.08)
                    .. controls (1.58, 1.25) and (1.52, 1.62) .. (1.12, 1.87)
                    .. controls (0.72, 2.18) and (0.63, 2.28) .. (0.5, 2.38) -- cycle;
                    
    \draw[gray!60, thick] \lineaPrediccion;
    
    \edef\ejecutarLoopDos{\noexpand\foreach \noexpand\x/\noexpand\y in {\puntosRegresion}{\noexpand\node[punto] at (\noexpand\x,\noexpand\y) {};}}
    \ejecutarLoopDos
\end{scope}

% --- GRÁFICO 3: Quantile Regression Score (CQR Asimétrico) ---
\begin{scope}[shift={(11,5)}]
    \node[font=\large\bfseries, text=black!70] at (2.2, 3.8) {Score CQR};
    \draw[eje] (0,0) -- (4.5,0) node[below left] {$\mathcal{X}$};
    \draw[eje] (0,0) -- (0,3.5) node[above left] {$\mathcal{Y}$};
    
    \fill[banda] 
        (0.5, 2.92) .. controls (0.79, 2.58) and (0.92, 2.42) .. (1.25, 2.22)
                    .. controls (2.12, 1.82) and (2, 2.75) .. (2.75, 3)
                    .. controls (3.71, 2.82) and (3.01, 1.69) .. (4.03, 1.1) --
        (4.03, 0.4) .. controls (3.38, 0.67) and (3.31, 1.08) .. (2.88, 1.42)
                    .. controls (2.45, 1.62) and (2.08, 1.48) .. (1.5, 1.62)
                    .. controls (1, 2.13) and (0.68, 2.28) .. (0.5, 2.38) -- cycle;

    \draw[gray!60, thick] \lineaPrediccion;
    
    \edef\ejecutarLoopTres{\noexpand\foreach \noexpand\x/\noexpand\y in {\puntosRegresion}{\noexpand\node[punto] at (\noexpand\x,\noexpand\y) {};}}
    \ejecutarLoopTres
\end{scope}

% =========================================================================
% FILA INFERIOR: CLASIFICACIÓN
% =========================================================================
\begin{scope}[shift={(84pt,-6pt)}]
    \draw[thick] (0,0) -- (9.5,0) node[below left] {$\mathcal{Y}$}; 
    \draw[thick] (0,0) -- (0,3.2); 
    
    \node[font=\LARGE] at (0.8, 3.3) {$\hat{\pi}(y \mid x)$};
    \node[font=\LARGE\bfseries] at (6.46, 3.3) {$\mathcal{C}(X_{n+1}) = \{1, 2, 3, 7\}$};
    
    \draw[barra_in] (0.4, 0) rectangle (0.9, 2.6);
    \node[below=2pt, black] at (0.65, 0) {$1$};
    
    \draw[barra_in] (1.3, 0) rectangle (1.8, 2.6);
    \node[below=2pt, black] at (1.55, 0) {$2$};
    
    \draw[barra_in] (2.2, 0) rectangle (2.7, 1.9);
    \node[below=2pt, black] at (2.45, 0) {$3$};
    
    \draw[barra_out] (3.1, 0) rectangle (3.6, 1.0);
    \draw[barra_out] (4.0, 0) rectangle (4.5, 1.1);
    \draw[barra_out] (4.9, 0) rectangle (5.4, 0.4);
    
    \draw[barra_in] (5.8, 0) rectangle (6.3, 1.5);
    \node[below=2pt, black] at (6.05, 0) {$7$};
    
    \draw[barra_out] (6.7, 0) rectangle (7.2, 0.8);
    \draw[barra_out] (7.6, 0) rectangle (8.1, 0.3);
    \draw[barra_out] (8.5, 0) rectangle (9.0, 0.3);
    
    \draw[thick, gray!90] (-0.3, 1.5) -- (9.8, 1.5);
    \node[right, font=\LARGE] at (9.9, 1.5) {$t^*(x)$};
\end{scope}
\end{tikzpicture}

\shorthandon{>}}
    \caption{Comparativa de métodos base de Predicción Conforme: enfoques continuos para regresión (arriba) y selección por umbral de probabilidad para clasificación (abajo).}
    \label{fig:conformal_full_panel}
\end{figure}

En la Figura~\ref{fig:conformal_full_panel} se resumen estas técnicas base de Predicción Conforme, a partir de las cuales se desarrollan formulaciones más elaboradas para adaptarse a las particularidades de cada dominio.



\section{Control de Riesgo Conforme}

En problemas con salidas estructuradas o de alta dimensión, como la clasificación multietiqueta, la detección de objetos o la predicción densa sobre grillas espaciales, exigir una cobertura marginal estricta sobre la totalidad de la salida simultáneamente suele resultar demasiado restrictivo y conduce a conjuntos de predicción poco informativos. Para superar esta limitación, el \textbf{Control de Riesgo Conforme} (\textit{Conformal Risk Control}, CRC) extiende el marco de la Predicción Conforme permitiendo controlar el valor esperado de una función de pérdida general $\ell$ en lugar de limitarse únicamente a la probabilidad de error binario. De hecho, CRC constituye una generalización directa de la Predicción Conforme: si se define la pérdida como la función indicadora de falta de cobertura $\ell(\mathcal{C}(X), Y) = \mathbb{I}\{Y \notin \mathcal{C}(X)\}$, la cota sobre el riesgo esperado $\mathbb{E}[\ell(\mathcal{C}(X), Y)] \leq \alpha$ recupera exactamente la garantía clásica de cobertura marginal $\mathbb{P}(Y \in \mathcal{C}(X)) \geq 1 - \alpha$.

Para operar dentro de este esquema general, se construye una familia de conjuntos de predicción $\mathcal{C}_\lambda(X)$ indexada por un parámetro escalar $\lambda$, de manera tal que al incrementar $\lambda$ los conjuntos se expandan de forma anidada y la pérdida asociada disminuya monotónicamente. Bajo el mismo supuesto de intercambiabilidad entre las muestras de calibración y el dato de prueba, CRC permite estimar un umbral óptimo $\hat{\lambda}$ a partir de un conjunto de calibración finito que garantiza matemáticamente que el riesgo esperado en futuras observaciones no exceda el nivel de tolerancia $\alpha \in (0,1)$ fijado por el usuario.

\subsection{Parametrización de conjuntos predictivos y clasificación LAC}

Para construir conjuntos de predicción parametrizados a partir de un modelo probabilístico previamente entrenado, se adopta la estrategia del clasificador basado en el conjunto menos ambiguo (\textit{Least Ambiguous Set-Valued Classifier}, LAC). Dado un predictor que produce estimaciones de probabilidad normalizadas (por ejemplo, salidas \textit{softmax}) denotadas por $\hat{f}_{ijk}(X)$, donde $(i,j)$ indexa la posición espacial o estructural dentro de la salida y $k$ representa cada clase del conjunto de etiquetas $\mathcal{Y} = \{1, \dots, K\}$, el conjunto de predicción conforme $\mathcal{C}_\lambda(X)_{ij}$ para cada posición se define incluyendo todas las clases cuya probabilidad estimada supere un umbral controlado por $\lambda \in [0,1]$:
\begin{equation}
    \mathcal{C}_\lambda(X)_{ij} = \{ k \in \mathcal{Y} : \hat{f}_{ijk}(X) \geq 1 - \lambda \}.
\end{equation}

El escalar $\lambda$ opera así como un mecanismo continuo de tolerancia. A medida que $\lambda$ aumenta, el umbral de corte $1-\lambda$ disminuye y el conjunto $\mathcal{C}_\lambda(X)_{ij}$ incorpora una mayor cantidad de clases candidatas para absorber la incertidumbre local del modelo. Esta construcción garantiza además la propiedad de anidamiento, es decir, si $\lambda \leq \lambda'$ entonces $\mathcal{C}_\lambda(X)_{ij} \subseteq \mathcal{C}_{\lambda'}(X)_{ij}$ para toda posición $(i,j)$, condición fundamental para asegurar que el error disminuya de forma monótona al aumentar el parámetro. En la práctica, cuando $\lambda$ toma valores muy pequeños, suele retenerse al menos la clase de mayor probabilidad estimada para evitar conjuntos vacíos.

\subsection{Calibración del umbral de riesgo}

Una vez definida la parametrización de los conjuntos $\mathcal{C}_\lambda(X)$, el objetivo de CRC consiste en determinar un valor $\hat{\lambda}$ tal que, ante una nueva observación de prueba $(X_{\text{test}}, Y_{\text{test}})$, el riesgo esperado bajo la función de pérdida elegida $\ell$ se mantenga acotado por el nivel $\alpha$:
\begin{equation}
    \mathbb{E}\big[ \ell(\mathcal{C}_{\hat{\lambda}}(X_{\text{test}}), Y_{\text{test}}) \big] \leq \alpha.
\end{equation}

Para estimar $\hat{\lambda}$ a partir de los datos, se evalúa primero el riesgo empírico $\hat{R}_n(\lambda)$ sobre un conjunto de calibración independiente compuesto por $n$ observaciones $\{(X_i, Y_i)\}_{i=1}^n$:
\begin{equation}
    \hat{R}_n(\lambda) = \frac{1}{n} \sum_{i=1}^n \ell\big(\mathcal{C}_\lambda(X_i), Y_i\big).
\end{equation}
Dado que $\hat{R}_n(\lambda)$ es una estimación empírica calculada sobre una muestra finita, elegir directamente el menor $\lambda$ que cumpla $\hat{R}_n(\lambda) \leq \alpha$ no basta para asegurar la cota teórica en datos nuevos. Asumiendo que la pérdida está acotada superiormente por una constante conocida $B$ (típicamente $B=1$) y es no creciente respecto de $\lambda$, se incorpora una corrección de muestra finita que define al estimador $\hat{\lambda}$ como el ínfimo del conjunto compatible:
\begin{equation}
    \hat{\lambda} = \inf \left\{ \lambda : \frac{n}{n+1} \hat{R}_n(\lambda) + \frac{B}{n+1} \leq \alpha \right\}.
\end{equation}

La dinámica de búsqueda de este umbral se ilustra en la Figura~\ref{fig:crc_optimization_flow}, donde el parámetro $\lambda$ cumple un rol análogo al del cuantil $\hat{q}$ en la Predicción Conforme por Partición.

\begin{figure}[H]
    \centering
    \resizebox{1.0\textwidth}{!}{\begin{tikzpicture}[scale=1.3]

    % --- EJES CARTESIANOS ---
    % Corregido el nodo para evitar conflictos con babel y caracteres activos
    \draw[-stealth, thick] (-0.2,0) -- (6.5,0) node[right] {$\lambda \in \Lambda$};
    \draw[-stealth, thick] (0,-0.2) -- (0,4.5) node[above] {Riesgo Empírico $\hat{R}_n(\lambda)$ [Costo]};
    
    % --- LÍNEAS DE REFERENCIA ---
    \draw[dashed, red, thick] (-0.2, 1.5) -- (6.0, 1.5) node[right, black] {$\alpha$ (Límite de Riesgo)};
    \draw[dotted, gray] (-0.2, 3.8) -- (6.0, 3.8) node[right] {$B=1$};
    \node[left] at (0, 3.8) {$1.0$};
    \node[left] at (0, 0) {$0.0$};

    % --- CURVA DE COSTO MONÓTONA ---
    \draw[ultra thick, blue, domain=0.2:5.8, samples=100] 
        plot (\x, {3.8 / (1 + exp(1.3*(\x - 2.5)))});
    \node[blue, right] at (4.5, 0.8) {$\hat{R}_n(\lambda) = \frac{1}{n}\sum_{i=1}^n L(Y_i, C_\lambda(X_i))$};

    % --- ESTADO A: REGIÓN DE INVIABILIDAD (Subcalibración) ---
    \coordinate (A) at (0.6, 3.5);
    \filldraw[black] (A) circle (2.5pt);
    \node[above right] at (A) {\textbf{Estado A}: $\lambda \to \lambda_{\min}$};
    % Corregido 'Bour' por la punta estándar 'stealth-'
    \draw[stealth-, thin, black!60!cyan] (A) -- (1.5, 4.2) node[right, font=\footnotesize, align=left, text=blue!50!black] {
        $|C_\lambda(X)| \to \min$ (Conjuntos vacíos)\\
        $\hat{R}_n(\lambda) \gg \alpha \implies \text{Región Inviable}$
    };

    % --- ESTADO B: TRÁNSITO GEOMÉTRICO (Monotonicidad) ---
    \coordinate (B) at (2.2, 2.3);
    \filldraw[black] (B) circle (2.5pt);
    \node[above right] at (B) {\textbf{Estado B}};
    % Corregido 'Bour' por la punta estándar 'stealth-'
    \draw[stealth-, thin, black!60!cyan] (B) -- (3.6, 3.2) node[right, font=\footnotesize, align=left, text=blue!50!green!80!black] {
        $\lambda \uparrow \implies C_{\lambda_{1}} \subseteq C_{\lambda_{2}}$\\
        $L(Y, C_{\lambda_{2}}) \leq L(Y, C_{\lambda_{1}})$\\
        El costo desciende continuamente
    };

    % --- ESTADO C: INTERSECCIÓN Y CONVERGENCIA ÓPTIMA ---
    \coordinate (C) at (2.9, 1.5); 
    \filldraw[red] (C) circle (3pt);
    \node[below left, red] at (C) {\textbf{Estado C}};
    
    % Proyección al eje lambda_hat
    \draw[dashed, thick, black] (2.9, 1.5) -- (2.9, 0) node[below, font=\large] {$\hat{\lambda}$};

    % Cuadro explicativo del Ínfimo Matemático
    \draw[black, fill=yellow!10, rounded corners, thin] (3.1, -1.5) rectangle (8.2, -0.4);
    \node[font=\small, align=left] (node1) at (5.7, -0.9) {
        \textbf{Optimización (Ínfimo Estadístico):}\\
        $\hat{\lambda} = \inf \left\{ \lambda : \frac{n}{n+1}\hat{R}_n(\lambda) + \frac{B}{n+1} \leq \alpha \right\}$
    };

    % --- ANOTACIÓN DE INFERENCIA EN TEST ---
    \node[draw, rounded corners, fill=blue!5, font=\footnotesize, align=left] at (0.9, -0.95) {
        \textbf{Testeo:}\\
        $C_{\hat{\lambda}}(X_{\text{test}}) = \{k : \hat{f}_k \geq 1 - \hat{\lambda}\}$\\
        $\mathbb{E}[L(Y_{\text{test}}, C_{\hat{\lambda}}(X_{\text{test}}))] \leq \alpha$
    };

\end{tikzpicture}}
    \caption{Dinámica de optimización del umbral $\hat{\lambda}$ en Control de Riesgo Conforme (CRC).}
    \label{fig:crc_optimization_flow}
\end{figure}

Tal como se observa en la Figura~\ref{fig:crc_optimization_flow}, partiendo de un valor inicial de $\lambda$ (Estado A), el riesgo empírico ajustado sobre el conjunto de calibración supera el umbral tolerado $\alpha$, por lo que el procedimiento incrementa $\lambda$ para ampliar los conjuntos predictivos. En una etapa intermedia (Estado B), aunque todavía no se alcanza la cota de seguridad, el riesgo desciende de manera consistente gracias a la propiedad de monotonía decreciente de la pérdida. Finalmente, el algoritmo converge al punto de corte $\hat{\lambda}$ (Estado C), donde la curva intersecta el nivel permitido por $\alpha$. Por la definición del ínfimo matemático, todo valor $\lambda \geq \hat{\lambda}$ satisface estrictamente la garantía de control de riesgo en muestra finita.

\subsection{Función de pérdida y métrica de eficiencia}

Para que la garantía de CRC sea válida, la función de pérdida de calibración $\ell$ debe estar acotada por $B$ y ser monótonamente no creciente a medida que los conjuntos de predicción se agrandan. En salidas estructuradas con múltiples posiciones, utilizar una pérdida estrictamente binaria que penalice con $1$ si al menos un único elemento queda fuera de su conjunto predicho resulta excesivamente severo, ya que obliga a seleccionar $\hat{\lambda}$ muy cercano a $1$ e incluir casi todas las clases en cada posición. En su lugar, resulta mucho más informativo emplear una penalización gradual denominada \textbf{pérdida por falta de cobertura}.

Para formularla, se representa la referencia verdadera mediante un tensor binario $Y \in \{0,1\}^{H \times W \times K}$ (en codificación \textit{one-hot}), donde $Y_{ijk} = 1$ si la clase real en la posición $(i,j)$ es $k$ y $0$ en caso contrario. De este modo, la suma total $\sum_{i,j,k} Y_{ijk}$ coincide exactamente con la cantidad de elementos válidos etiquetados $n_{\text{pixels}}$. De forma análoga, la predicción conforme $\mathcal{C}_\lambda(X)$ induce un tensor binario $Z \in \{0,1\}^{H \times W \times K}$ definido por $Z_{ijk} = \mathbb{I}(k \in \mathcal{C}_\lambda(X)_{ij})$, que toma el valor $1$ si la clase $k$ fue incluida en el conjunto predicho para la posición $(i,j)$ y $0$ en caso contrario.

Bajo esta notación, el producto $Z_{ijk} Y_{ijk}$ vale $1$ si y solo si la clase verdadera $k$ de la posición $(i,j)$ quedó contenida dentro del conjunto de predicción $\mathcal{C}_\lambda(X)_{ij}$. Al sumar estos aciertos sobre toda la estructura y dividir por el total de elementos válidos $\sum_{i,j,k} Y_{ijk}$, se obtiene la fracción empírica de posiciones correctamente cubiertas. Su complemento respecto de la unidad define entonces la proporción de elementos no cubiertos:
\begin{equation}
    \ell(Z, Y) = 1 - \frac{\sum_{i,j,k} Z_{ijk} Y_{ijk}}{\sum_{i,j,k} Y_{ijk}}.
\end{equation}
Esta función toma valores en el intervalo $[0,1]$, por lo que verifica la cota superior $B = 1$, y disminuye de forma monótona al aumentar $\lambda$.

Por otra parte, dado que la pérdida $\ell(Z,Y)$ podría reducirse trivialmente a cero incluyendo todas las $K$ clases en cada posición ($Z_{ijk} = 1$ para todo $k$), es necesario acompañar el control de riesgo con una métrica que evalúe la \textbf{eficiencia} de los conjuntos generados. Para una posición individual $(i,j)$, la suma $\sum_{k=1}^K Z_{ijk}$ contabiliza cuántas clases candidatas fueron activadas simultáneamente dentro de $\mathcal{C}_\lambda(X)_{ij}$. Al promediar esta cantidad sobre el total de elementos válidos $n_{\text{pixels}}$, se obtiene la \textbf{tasa de activación} ($\text{AR}$):
\begin{equation}
    \text{AR}(Z) = \frac{1}{n_{\text{pixels}}} \sum_{i,j,k} Z_{ijk}.
\end{equation}
Esta métrica cuantifica el tamaño medio de los conjuntos de predicción por posición. Un valor de $\text{AR}(Z)$ cercano a $1$ indica alta confianza del modelo, ya que el conjunto predicho se reduce prácticamente a una única clase por elemento, mientras que valores mayores reflejan un aumento en la incertidumbre local que obliga al método a retener múltiples etiquetas candidatas para sostener el nivel de riesgo $\alpha$ exigido. En el capítulo experimental, esta formulación general se aplicará directamente a tareas de segmentación semántica a nivel de píxel.


























